\documentclass[12pt]{article}

\usepackage{fontspec}
\newfontfamily\handwriting{a-casual-handwritten-pen}[Path=../assets/fonts/,Extension=.ttf,Scale=1.5]
\newfontfamily\comic{Comic Neue}

\usepackage[utf8]{inputenc}
\usepackage{amsmath}
\usepackage{amssymb}
\usepackage{enumitem}
\usepackage{xcolor}
\usepackage{soul}
\usepackage{tikz}
\usepackage[most]{tcolorbox}
\usepackage{titlesec}
\usepackage{multicol}
\usepackage{environ}

\usetikzlibrary{fit, positioning, arrows.meta}

% --- Custom Colors ---
\definecolor{primaryblue}{RGB}{42, 111, 151}
\definecolor{exbg}{RGB}{255, 240, 245}
\definecolor{rulebg}{RGB}{232, 245, 233}
\definecolor{highlight}{RGB}{255, 243, 176}
\definecolor{pinkanno}{RGB}{255, 105, 180}
\definecolor{purpleanno}{RGB}{147, 112, 219}
\definecolor{solutionbg}{RGB}{245, 245, 255}

% --- Header Styling ---
\titleformat{\section}{\color{primaryblue}\normalfont\Large\bfseries}{\thesection}{1em}{}
\titleformat{\subsection}{\color{primaryblue}\normalfont\large\bfseries}{\thesubsection}{1em}{}

% --- Friendly Boxes ---
% --- Auto-numbered Example Box ---
\newcounter{examplenum}

\newtcolorbox{friendlyexbox}[1]{
    colback=exbg,
    colframe=pinkanno!50,
    arc=5pt,
    title=\textbf{#1},
    coltitle=black,
    attach title to upper,
    after title={:\quad}
}

\newenvironment{friendlyex}{%
    \refstepcounter{examplenum}%
    \begin{friendlyexbox}{Example \theexamplenum}%
}{%
    \end{friendlyexbox}%
}

\newtcolorbox{friendlyrule}[1]{
    colback=rulebg,
    colframe=primaryblue!30,
    arc=5pt,
    fonttitle=\bfseries,
    title={#1},
    coltitle=primaryblue
}

\newcommand{\funhl}[1]{\sethlcolor{highlight}\hl{#1}}

% --- Show/Hide Solutions ---
\newif\ifshowsolutions

% Uncomment the next line to show solutions.
\showsolutionstrue

\newsavebox{\solutionmeasurebox}

\NewEnviron{solutionbox}{
\ifshowsolutions
\begin{tcolorbox}[
    colback=solutionbg,
    colframe=purpleanno!45,
    arc=5pt,
    title={Solution},
    coltitle=primaryblue,
    fonttitle=\bfseries
]
\BODY
\end{tcolorbox}
\else
\sbox{\solutionmeasurebox}{\begin{minipage}{\dimexpr\linewidth-10pt\relax}\BODY\end{minipage}}
\vspace{\dimexpr\ht\solutionmeasurebox+\dp\solutionmeasurebox+3em\relax}
\fi
}

\begin{document}
\handwriting

\begin{center}
    \Large\textbf{\color{primaryblue}Module 1 Lesson 6\\Absolute Value Equations and Inequalities} \\
    \vspace{0.2cm}
    \hrule
\end{center}

\section*{Objectives}

\begin{friendlyrule}{In this lesson, we will learn how to}
\begin{itemize}
    \item define absolute value as distance from zero;
    \item solve absolute value equations;
    \item solve absolute value inequalities;
    \item apply absolute value to real-world tolerance problems.
\end{itemize}
\end{friendlyrule}

\section*{1. Absolute Value}

\begin{friendlyrule}{Absolute Value}
The \funhl{absolute value} of a real number $x$ is its distance from zero, denoted $|x|$. For example,

\[
|5|=5,\qquad |-5|=5.
\]
\end{friendlyrule}

\section*{2. Absolute Value Equations}

\begin{friendlyrule}{Solving $|X|=a$}
For $a>0$ and an algebraic expression $X$,

\[
|X|=a\quad\text{is equivalent to}\quad X=-a\ \ \text{or}\ \ X=a.
\]

For example, if $|x|=5$, then $x=-5$ or $x=5$.
\end{friendlyrule}

\begin{friendlyex}
Solve each absolute value equation.

\begin{itemize}
    \item[(a)] $8|x-3|=4$
    \item[(b)] $|3x-6|+5=3$
    \item[(c)] $5-|x+2|=3$
\end{itemize}
\end{friendlyex}

 

\begin{solutionbox}
\textbf{(a)} Divide both sides by $8$:

\[
|x-3|=\frac12.
\]

\[
x-3=\frac12\quad\text{or}\quad x-3=-\frac12.
\]

\[
x=\frac72\quad\text{or}\quad x=\frac52.
\]

\textbf{(b)} Isolate the absolute value:

\[
|3x-6|=3-5=-2.
\]

An absolute value can never equal a negative number, so there is \textbf{no solution}.

\textbf{(c)} Isolate the absolute value:

\[
-|x+2|=3-5=-2\ \Longrightarrow\ |x+2|=2.
\]

\[
x+2=2\quad\text{or}\quad x+2=-2.
\]

\[
x=0\quad\text{or}\quad x=-4.
\]
\end{solutionbox}

  

\section*{3. Absolute Value Inequalities}

\begin{friendlyrule}{Solving Absolute Value Inequalities}
For $a>0$ and an algebraic expression $X$:

\[
|X|<a\quad\text{is equivalent to}\quad -a<X<a.
\]

\[
|X|>a\quad\text{is equivalent to}\quad X<-a\ \ \text{or}\ \ X>a.
\]
\end{friendlyrule}

\begin{friendlyex}
Solve, and write the solution in interval notation.

\begin{itemize}
    \item[(a)] $|4x|<12$
    \item[(b)] $|5x+1|<8$
    \item[(c)] $|5x+1|>8$
\end{itemize}
\end{friendlyex}

 

\begin{solutionbox}
\textbf{(a)}

\[
-12<4x<12.
\]

\[
-3<x<3.
\]

Interval notation: $(-3,3)$.

\textbf{(b)}

\[
-8<5x+1<8.
\]

\[
-9<5x<7.
\]

\[
-\frac95<x<\frac75.
\]

Interval notation: $\left(-\dfrac95,\dfrac75\right)$.

\textbf{(c)}

\[
5x+1<-8\quad\text{or}\quad 5x+1>8.
\]

\[
5x<-9\quad\text{or}\quad 5x>7.
\]

\[
x<-\frac95\quad\text{or}\quad x>\frac75.
\]

Interval notation: $\left(-\infty,-\dfrac95\right)\cup\left(\dfrac75,\infty\right)$.
\end{solutionbox}

  

\section*{4. Applications}

\begin{friendlyex}
\begin{itemize}
    \item[(a)] Suppose you want your room temperature to stay within $4$ degrees of $70^\circ$F. Determine the temperature range using an absolute value inequality.
    \item[(b)] A machine is designed to produce metal rods that are $10$ cm long, with a maximum allowable error of $0.03$ cm. Write and solve an absolute value inequality to determine the acceptable rod lengths.
\end{itemize}
\end{friendlyex}

 

\begin{solutionbox}
\textbf{(a)} Let $T$ be the temperature. We need

\[
|T-70|\le 4.
\]

\[
-4\le T-70\le 4.
\]

\[
66\le T\le 74.
\]

The acceptable temperature range is $66^\circ$F to $74^\circ$F.

\textbf{(b)} Let $L$ be the length of the rod. We need

\[
|L-10|\le 0.03.
\]

\[
-0.03\le L-10\le 0.03.
\]

\[
9.97\le L\le 10.03.
\]

The acceptable rod lengths are between $9.97$ cm and $10.03$ cm.
\end{solutionbox}

  

\section*{Summary}

\begin{friendlyrule}{Key Ideas}
\begin{itemize}
    \item $|x|$ measures the distance of $x$ from $0$ on the number line, and is always $\ge 0$.
    \item $|X|=a$ (for $a>0$) splits into two cases: $X=a$ or $X=-a$.
    \item If isolating the absolute value produces a negative number on the other side, the equation has no solution.
    \item $|X|<a$ becomes a ``between'' (double) inequality: $-a<X<a$.
    \item $|X|>a$ splits into an ``or'' compound inequality: $X<-a$ or $X>a$.
    \item Absolute value inequalities are natural for describing tolerances and acceptable ranges.
\end{itemize}
\end{friendlyrule}

\end{document}
